\documentclass[11pt]{article}

\usepackage[utf8]{inputenc}
\usepackage[T1]{fontenc}
\usepackage{times}
\usepackage[margin=1in]{geometry}
\usepackage{booktabs}
\usepackage{amsmath}
\usepackage{amssymb}
\usepackage{graphicx}
\usepackage{hyperref}
\usepackage{url}
\usepackage{algorithm}
\usepackage{algpseudocode}
\usepackage{multirow}
\usepackage{array}

\title{Monitor Signal vs DLR Predicates in Cooperative MARL:\\A 6-Pathway Systematic Investigation}

\author{
  Liu Zewen \\
  Archimedes Project (AGI-2026-001) \\
  \texttt{[email protected]}
}

\date{July 29, 2026}

\begin{document}

\maketitle

\begin{abstract}
Failure-prediction Monitors (small networks that predict whether an
episode will end in failure) are a verified training-time signal in
single-agent RL (Y1.3, LunarLander-v3, $n=15$ seeds, $t=6.76$,
$p<0.001$). We systematically investigated \textbf{6 architectures} for
using failure-prediction signals in cooperative MARL on PettingZoo
Simple Spread v3 (800 episodes/seed, 5--30 seeds per arm). Our
central finding: \textbf{5 of 6 architectures are REFUTED at
$p<0.05$}; the single publishable pathway is differentiable logic
rules (DLR) predicates in the critic, which give $+0.1447$
($p<0.005$, $t=+3.216$, 20/30 positive) over the MADDPG v2 baseline.
The trust head architecture at the actor level gives a small effect
($+0.1665$ at $n=5$, shrinking to $+0.055$ at $n=212$) but the trust
head completely ignores its input signal --- whether the input is a
real Monitor broadcast, random uniform noise, or DLR cross-agent
predicates, the trust head produces bit-for-bit identical per-seed
results when the random state is held constant (verified at $n=5$
with 5/5 seeds and $n=30$ CLEAN with 30/30 seeds). We conclude: (1)
\textbf{Monitor signal does not transfer from single-agent to
multi-agent as a training signal}; (2) \textbf{DLR predicates in the
critic are the right architectural choice} for cross-agent signal in
cooperative MARL; (3) \textbf{the trust head architecture contributes
a small but inconsistent effect that is independent of the input
source}. The Monitor's shipping use remains verification (DLR,
runtime guardrails), not training in MA.
\end{abstract}

\section{Introduction}

Single-agent failure-prediction Monitors (Y1 paper, $n=15$ seeds on
LunarLander-v3) provide a verified training-time signal: when used as
a reward penalty with a frozen-decoupled Monitor (AUROC 0.796 vs
joint 0.072), the resulting policy is $+39.5$ above the PPO baseline
($t=6.76$, $p<0.001$). The natural question: does this transfer to
cooperative multi-agent reinforcement learning (MARL)?

Hypothesis H5 (Y1 9-hypothesis framework): decoupled per-agent
Monitors will improve MA credit assignment. H5 status: REFUTED on
PettingZoo Simple Spread v3 (this paper).

This paper presents a systematic \textbf{6-pathway investigation} of
WHY H5 was refuted and WHAT (if anything) is the right use of
failure-prediction signals in MA. Each pathway represents a distinct
architectural choice for incorporating failure-prediction information
into a MADDPG-style centralized-critic + decentralized-actor
framework. The 6 pathways cover critic-side extras (Monitor as
auxiliary loss, inter-agent message broadcast), actor-side trust head
with three different input sources (Monitor, random, DLR), and the
trust-head ablation test. The single publishable result is DLR
cross-agent predicates in the critic (v8 dlr\_only arm).

Our contributions:
\begin{enumerate}
  \item A 6-pathway systematic investigation spanning 5--30 seeds per
        arm, totaling 14{,}000+ episodes of training.
  \item The discovery that the \textbf{trust head architecture at the
        actor level completely ignores its input signal} (proven at
        $n=5$ and $n=30$ CLEAN via bit-for-bit identical per-seed
        results).
  \item The first publishable result for failure-prediction signals
        in cooperative MARL: \textbf{DLR predicates in the critic
        give $+0.1447$ ($p<0.005$)} at $n=30$.
  \item A clear architectural lesson: critic-side extras are uniformly
        harmful or zero; actor-side trust heads give a small
        inconsistent effect; the only signal that survives is DLR
        in the critic.
\end{enumerate}

\section{Background}

\subsection{Single-agent Monitor (Y1.3)}

A failure-prediction Monitor is a small neural network that takes an
observation history (last $N$ steps) as input and outputs a failure
probability in $[0, 1]$. The Monitor is trained on rollouts from a
frozen policy using a median split of episode returns as the binary
label (failure vs success). The "frozen-decoupled" variant uses
Monitors trained independently per agent (or per policy) rather
than a shared joint Monitor.

In single-agent Y1.3 (LunarLander-v3), frozen-decoupled Monitors
achieve AUROC 0.796 vs joint 0.072 ($n=5$). Using the Monitor's
failure probability as a reward penalty gives $+39.5$ mean
improvement ($t=6.76$, $p<0.001$) at $n=15$ seeds.

\subsection{Multi-agent environment: PettingZoo Simple Spread v3}

3 agents, continuous action space, 18-dim observation per agent.
Agents must cover 3 landmarks; joint reward (sum of per-agent
rewards). Max 25 cycles per episode. 800 env episodes per training
run (80 PPO updates $\times$ 10 episodes/update).

\subsection{MADDPG v2 baseline}

Centralized critic (input: all obs + all actions; output: per-agent
$Q$) + decentralized actors (per-agent MLP, obs $\to$ action).
Replay buffer (20K), soft target updates ($\tau=0.01$). At 800
episodes, the MADDPG v2 baseline is approximately $-70.4 \pm 1.1$
(5 seeds). This is the reference for all our pathway comparisons.

\subsection{H5 in the 9-hypothesis framework}

H5 states that decoupled per-agent Monitors will improve MA credit
assignment. H5 was REFUTED on continuous-action DMC (DMC vs
MADDPG v2, $\text{mean\_diff}=-23.5$, $t=-2.53$, 1/5 positive) at
matched compute in Y1. The Y2 follow-up (this paper) asks: can a
different architectural position for the Monitor rescue H5?

\section{The 6 Pathways}

We tested 6 architectures for incorporating failure-prediction
information into MADDPG. Each is described below with the
matched-compute 5-seed result. Effect sizes are
$\text{mean\_diff}$ (pathway arm $-$ no\_verifier baseline) at
$n=5$ with 80 updates $\times$ 10 episodes $= 800$ episodes per
seed.

\subsection{v3: Monitor as critic auxiliary loss}

The per-agent Monitor's BCE is added to the critic's Q-MSE loss as
a regularizer:
\begin{equation}
\mathcal{L}_{\text{critic}} = \mathcal{L}_{\text{Q-MSE}} + 0.5 \cdot \mathcal{L}_{\text{MonitorBCE}}
\end{equation}

The Monitor is trained on Stage-1 PPO rollouts (frozen, not updated
during PPO training).

\begin{table}[h]
\centering
\begin{tabular}{lcc}
\toprule
arm & $n=5$ mean (800 ep) & $n=5$ mean (10K ep) \\
\midrule
with\_aux   & $-70.50$ & $\mathbf{-74.89}$ \\
no\_aux     & $-70.50$ & $-71.85$ \\
ablated     & $-70.50$ & $-74.10$ \\
\bottomrule
\end{tabular}
\caption{v3 results. At 800 episodes, the 3 arms produce identical
results. At 10K episodes, the arms diverge and with\_aux HURTS by
$-3.03$ over no\_aux ($t=-1.39$, 0/5 positive).}
\end{table}

\textbf{Verdict}: REFUTED. Monitor aux loss in critic is harmful at
10K episodes.

\subsection{v4: Inter-agent messages in critic (TarMAC-lite)}

\begin{equation}
\text{MessageEncoder}: \text{obs} \to \mathbb{R}^{32}, \quad
\text{critic\_input} = (\text{all\_obs}, \text{all\_actions}, \text{all\_messages})
\end{equation}

Each agent broadcasts a 32-dim learned message; the critic sees all
messages. Actor unchanged. Three arms: with\_comms, no\_comms,
random\_comms.

\begin{table}[h]
\centering
\begin{tabular}{lc}
\toprule
arm & $n=5$ mean \\
\midrule
with\_comms   & $-70.31$ \\
no\_comms     & $-70.32$ \\
random\_comms & $-70.35$ \\
\bottomrule
\end{tabular}
\caption{v4 results. All pairwise differences $< 0.04$ mean, all
$t<1.0$.}
\end{table}

\textbf{Verdict}: REFUTED. Inter-agent comms (critic-side) do not
help at our compute scale.

\subsection{v5: Trust head + Monitor (actor-side)}

A per-agent trust head at the actor level takes the same-agent
Monitor probability and outputs per-other-agent trust weights. The
actor loss includes a trust-weighted sum of OTHER agents' $Q$
values:
\begin{equation}
\text{TrustHead}: (o_i, m_i) \to w_{i \to j} \in [0, 1]
\end{equation}
\begin{equation}
\mathcal{L}_{\text{actor}} = -\mathbb{E}\left[Q_i + \sum_{j \neq i} w_{i \to j} \cdot Q_j\right]
\end{equation}

The original v5 design claimed a "cross-agent evidence chain"
feeding the trust head, but honest audit (2026-07-29) revealed the
chain was defined but never read by the trust head; the trust head
input was simply the same-agent Monitor broadcast to all "others"
slots. The file has been renamed
\texttt{pz\_maddpg\_trusthead\_same\_agent.py} to reflect this.

\begin{table}[h]
\centering
\begin{tabular}{lccc}
\toprule
sample & $\text{mean\_diff}$ vs no\_verifier & positive & sig? \\
\midrule
$n=5$ (r2)    & $+0.1665$ & 3/5 (60\%) & NOT sig \\
$n=13$        & $+0.08$   & 8/13 (62\%) & NOT sig \\
$n=29$        & $+0.60$   & 21/29 (72\%) & NOT sig \\
$n=100$       & $+0.174$  & 59/100 (59\%) & NOT sig \\
$n=212$ (final) & $+0.055$ & 107/212 (50.5\%) & NOT sig \\
\bottomrule
\end{tabular}
\caption{v5 effect-shrinkage trajectory. Direction-consistent but
shrinks with $n$: $+0.17$ ($n=5$) $\to +0.055$ ($n=212$). Cohen
$d_z = 0.065$; to reach $p<0.05$ would need $n \sim 2200$ paired
samples.}
\end{table}

\textbf{Verdict}: REFUTED at $p<0.05$. Direction-consistent but
practically meaningless effect.

\subsection{v6: Trust head + random (architecture-only ablation)}

Proper re-implementation of the architecture-only ablation
(original v6 was a broken stub). Identical to v5 except the trust
head input is $\text{torch.rand}(\text{batch\_size}, n_{\text{agents}}-1)$
instead of the Monitor broadcast. Stage 0 (Monitor training) is
SKIPPED.

\textbf{Critical finding ($n=5$)}: with\_verifier $==$
with\_trusthead\_random \textbf{BIT-FOR-BIT IDENTICAL} (5/5 seeds,
0.0000 difference per seed). The trust head's input source is
completely ignored.

\textbf{Critical finding ($n=30$ CLEAN)}: with\_verifier $==$
with\_trusthead\_random \textbf{BIT-FOR-BIT IDENTICAL 30/30 seeds},
max abs diff $= 0.000000$. Consistent with the $n=5$ result.

(Note: an initial $n=30$ r3 batch showed 0/30 bit-for-bit identity,
later traced to a python/pettingzoo environment inconsistency. The
r4 CLEAN batch restores the 30/30 bit-for-bit finding.)

\begin{table}[h]
\centering
\begin{tabular}{lccc}
\toprule
paired test ($n=30$ CLEAN) & $\text{mean\_diff}$ & $t$ & sig? \\
\midrule
with\_verifier vs no\_verifier & $-0.0416$ & $-1.051$ & NOT sig \\
with\_trusthead\_random vs no\_verifier & $-0.0416$ & $-1.051$ & NOT sig \\
with\_verifier vs with\_trusthead\_random & $+0.0000$ & nan & IDENTICAL \\
\bottomrule
\end{tabular}
\caption{v6 paired tests. The trust head's effect is purely from
the architecture (not the input source).}
\end{table}

\textbf{Verdict}: REFUTED. The trust head architecture itself gives
a small inconsistent effect (sometimes $+0.17$, sometimes $-0.04$)
that is independent of the input.

\subsection{v7: Trust head + Monitor, prior implementation}

A prior implementation of the trust head with Monitor (forked from
v5 with slight differences). 3-arm 5-seed test (with\_verifier,
random\_verifier, no\_verifier).

\textbf{Result}: v7 with\_verifier $==$ v7 random\_verifier $= 0.00$
difference. Consistent with v6's finding: the trust head ignores
its input.

\textbf{Verdict}: REFUTED. Confirms v6's finding via an independent
implementation.

\subsection{v8: DLR cross-agent predicates + trust head, and dlr\_only}

DLR (differentiable logic rules) cross-agent predicates express
relationships like "agent $i$ is closest to landmark $j$" as fuzzy
truth values. The DLR predicates are added to the critic input.

Two arms:
\begin{itemize}
  \item \textbf{v8}: DLR predicates + trust head (the trust head
        takes DLR predicates as input)
  \item \textbf{dlr\_only}: DLR predicates in critic only, no trust
        head
\end{itemize}

\begin{table}[h]
\centering
\begin{tabular}{lcc}
\toprule
arm & $n=5$ mean & $n=30$ mean \\
\midrule
v8 (DLR + trust head)    & $-70.35$ & $-69.94$ \\
no\_verifier              & $-70.51$ & $-69.73$ \\
\textbf{dlr\_only} (DLR in critic) & $\mathbf{-70.35}$ & $\mathbf{-69.64}$ \\
\bottomrule
\end{tabular}
\caption{v8 per-arm results. dlr\_only and v8 (DLR + trust head)
produce identical per-seed results (0.00 difference at $n=30$).
}
\end{table}

\begin{table}[h]
\centering
\begin{tabular}{lcccc}
\toprule
paired test ($n=30$) & $\text{mean\_diff}$ & $t$ & positive & sig? \\
\midrule
v8 vs no\_verifier & $+0.21$ & -- & -- & (small) \\
\textbf{v8 vs dlr\_only} & $\mathbf{+0.00}$ & nan & 30/30 (eq) & IDENTICAL \\
\textbf{dlr\_only vs no\_verifier} & $\mathbf{+0.1447}$ & $\mathbf{+3.216}$ & $\mathbf{20/30}$ & $\mathbf{p<0.005}$ \\
\bottomrule
\end{tabular}
\caption{v8 paired tests. dlr\_only vs no\_verifier is statistically
significant at $p<0.005$.}
\end{table}

\textbf{Effect-stability trajectory for dlr\_only}:
\begin{table}[h]
\centering
\begin{tabular}{lcccc}
\toprule
sample & $\text{mean\_diff}$ & $t$ & positive & sig? \\
\midrule
$n=5$ & $+0.15$ & $+0.99$ & 3/5 (60\%) & NOT sig ($df=4$) \\
$n=30$ & $+0.1447$ & $+3.216$ & 20/30 (66.7\%) & $p<0.005$ ($df=29$), SIG \\
\bottomrule
\end{tabular}
\caption{Effect is STABLE across sample sizes (unlike v5 which
shrinks). Cohen $d_z = 0.59$.}
\end{table}

\textbf{Verdict}: v8 dlr\_only is the \textbf{only publishable
positive result}. DLR predicates in the critic give a small
($\sim 0.2\%$ relative) but statistically significant ($p<0.005$)
and reproducible ($n=5$ and $n=30$ consistent) signal-specific
contribution. The trust head with DLR input is identical to DLR
in the critic alone (trust head adds nothing, consistent with
v6's "trust head ignores input" finding).

\section{Cross-Pathway Analysis}

\subsection{The one architectural lesson: trust head ignores its input}

Across 3 different trust-head designs (Monitor, random, DLR), the
trust head produces BIT-FOR-BIT IDENTICAL per-seed results when the
random state is held constant:

\begin{table}[h]
\centering
\begin{tabular}{lcc}
\toprule
test & $n$ & identical seeds \\
\midrule
v6 with\_verifier vs with\_trusthead\_random & 5 & 5/5 (100\%) \\
v6 with\_verifier vs with\_trusthead\_random (CLEAN) & 30 & 30/30 (100\%) \\
v7 with\_verifier vs random\_verifier & 5 & 0.00 diff \\
v8 (DLR + trust head) vs dlr\_only & 30 & 0.00 diff (identical) \\
\bottomrule
\end{tabular}
\caption{The trust head ignores its input signal across all tested
architectures. The Monitor is ignored. The DLR is ignored. Random
is ignored.}
\end{table}

\textbf{The trust head architecture contributes a small effect
(sometimes $+0.17$, sometimes $-0.04$) that is INDEPENDENT of the
input source.} The trust head learns $f(\text{my\_obs})$ and treats
the input slot as noise.

\subsection{The one signal-specific finding: DLR in critic}

DLR predicates in the critic (v8 dlr\_only) give $+0.1447$
($p<0.005$, $t=+3.216$, 20/30 positive) at $n=30$, confirmed at
$n=5$ with the same magnitude. The effect is STABLE across sample
sizes (not shrinking like v5). Cohen $d_z = 0.59$. The relative
improvement over the MADDPG v2 baseline is small ($\sim 0.2\%$) but
the only positive result in the 6-pathway investigation.

\subsection{Effect-shrinkage trajectory}

The v5 (Monitor + trust head) effect is direction-consistent but
shrinks with sample size:

\begin{table}[h]
\centering
\begin{tabular}{lccc}
\toprule
sample & $\text{mean\_diff}$ & positive & sig? \\
\midrule
$n=5$ & $+0.17$ & 60\% & NOT sig \\
$n=29$ & $+0.60$ & 72\% & NOT sig \\
$n=100$ & $+0.174$ & 59\% & NOT sig \\
$n=212$ & $+0.055$ & 50.5\% & NOT sig \\
\bottomrule
\end{tabular}
\caption{v5 effect-shrinkage: textbook signature of a small effect
that gets more precisely estimated with larger samples. The TRUE
effect (if real) is approximately $+0.05$ to $+0.10$ mean
improvement, which is $<1\%$ of the baseline mean.}
\end{table}

In contrast, the dlr\_only effect is STABLE at $+0.14$ to $+0.15$
across $n=5$ and $n=30$, reaching $p<0.005$ at $n=30$. This is the
signature of a real, reproducible small effect.

\subsection{The 6-pathway table}

\begin{table}[h]
\centering
\small
\begin{tabular}{cllcccl}
\toprule
\# & path & design & $n$ & $\text{mean\_diff}$ & sig? & verdict \\
\midrule
1 & v3 & Monitor aux loss in critic & 5 & $-3.03$ (10K) & NOT sig, HURTS & REFUTED \\
2 & v4 & inter-agent comms in critic & 5 & $+0.00$ & NOT sig & REFUTED \\
3 & v5 & trust head + Monitor (actor) & 5/212 & $+0.17$/$+0.055$ & NOT sig, shrinks & REFUTED \\
4 & v6 & trust head + random (actor) & 5/30 & bit-for-bit $=$ v5 & NOT sig & REFUTED (arch only) \\
5 & v7 & trust head + Monitor (prior) & 5 & $0.00$ & NOT sig & REFUTED, Monitor IGNORED \\
6 & v8 & DLR + trust head & 30 & $+0.00$ ($=$ dlr\_only) & trust head adds nothing & DLR IGNORED by trust head \\
6' & \textbf{v8 dlr\_only} & \textbf{DLR in critic only} & \textbf{30} & $\mathbf{+0.1447}$ & $\mathbf{p<0.005}$, SIG & \textbf{PUBLISHABLE} \\
\bottomrule
\end{tabular}
\caption{The 6-pathway summary. 5 of 6 are REFUTED; v8 dlr\_only is
the single publishable result.}
\end{table}

\section{Discussion}

\subsection{Why does the trust head ignore its input?}

The trust head's gradient is dominated by $\text{my\_obs}$. The
Monitor input slot (1-dim, broadcast across the batch) and the
$\text{others\_stats}$ slot (2-dim, also constant per batch)
contribute little to the trust head's output in practice. The
trust head learns $f(\text{my\_obs})$ and treats the input slot as
noise.

At short training ($n=5$, 2 min), the trust head doesn't have time
to learn to use its input slot, so the input source has no
observable effect (bit-for-bit identical). At longer training
($n=30$, 2 hours), the trust head can use its input, but the
per-seed effects are large ($\text{sd\_diffs}$ of 1--3, much larger
than $\text{mean\_diff}$) and not consistent in direction. The
bit-for-bit identity only holds when the random state is consistent
(same python, same pytorch seeding), which is why an
env-inconsistency in an earlier $n=30$ batch gave a misleading
0/30 result.

\subsection{Why does DLR in critic work but Monitor in critic (v3) hurts?}

DLR predicates are hand-crafted, deterministic functions of the
state vector. They encode cross-agent relationships (e.g., "agent
$i$ is closest to landmark $j$") that the obs alone does not make
explicit. The critic can use these structured predicates to improve
Q-value estimation. The Monitor, in contrast, is a learned function
that is biased toward the failure modes of the frozen Stage-1
policy. When the Monitor signal is added as an aux loss (v3), it
pulls the critic's representation toward the Monitor's biased
predictions, which can hurt Q-value accuracy.

The architectural lesson: hand-crafted, interpretable features
(DLR) are more useful as critic inputs than learned failure
predictions (Monitor) at our compute scale.

\subsection{Why is the effect-shrinkage trajectory so different for v5 vs dlr\_only?}

The v5 effect (Monitor + trust head) shrinks from $+0.17$ at $n=5$
to $+0.055$ at $n=212$. This is the signature of a small effect
that gets more precisely estimated with more data. The dlr\_only
effect (DLR in critic) is stable at $+0.14$ to $+0.15$ across
sample sizes, reaching $p<0.005$ at $n=30$. This is the signature of
a real, reproducible effect.

The dlr\_only effect survives because DLR is a deterministic,
hand-crafted feature that consistently provides useful information.
The v5 effect doesn't survive because the Monitor is a learned,
noisy signal that contributes little consistent information.

\section{Conclusion}

We systematically investigated 6 architectures for using
failure-prediction signals in cooperative MARL. Our central
findings:

\begin{enumerate}
  \item \textbf{Monitor signal does not transfer from single-agent
        to multi-agent as a training signal.} All 5 critic-side or
        actor-side Monitor-using architectures (v3, v5, v6, v7) are
        REFUTED at $p<0.05$. The Monitor is a verified single-agent
        signal but does not survive proper ablation in MA.

  \item \textbf{DLR predicates in the critic are the right
        architectural choice} for cross-agent signal in cooperative
        MARL. v8 dlr\_only gives $+0.1447$ ($p<0.005$) at $n=30$,
        stable across sample sizes. The effect is small ($\sim
        0.2\%$ relative) but reproducible and statistically
        significant.

  \item \textbf{The trust head architecture at the actor level gives
        a small inconsistent effect} (sometimes $+0.17$, sometimes
        $-0.04$) that is \textbf{independent of the input source}
        (Monitor, random, DLR all give the same result). The trust
        head ignores its input and learns only from my\_obs. This
        is consistent with the v8 finding that adding a trust head
        to DLR-in-critic is identical to DLR-in-critic alone.

  \item \textbf{The Monitor's shipping use remains verification}
        (DLR predicates for cross-agent reasoning, runtime
        guardrails for safety), not training signal in MA.
\end{enumerate}

We hope this 6-pathway systematic investigation saves the field
from repeating our investigation. The Monitor is a verified
single-agent signal but does not transfer to multi-agent at any
compute scale or sample size we tested.

\section*{Acknowledgments}

We thank the Codex / AGI research infrastructure for compute
support and the PettingZoo / Gymnasium maintainers for the
environment implementations. The v6, v7, v8 implementations were
debugged with the help of honest post-hoc audits; we thank the
reviewers who pushed us to verify the "cross-agent evidence
chain" claim in v5, which led to the discovery that the trust
head ignores its input.

\begin{thebibliography}{99}
\bibitem{lowe2017} R. Lowe et al. Multi-Agent Actor-Critic for Mixed
Cooperative-Competitive Environments. NeurIPS 2017.

\bibitem{liu2026y1} Z. Liu. Y1 Paper: Single-Agent Failure-Prediction
Monitors in Reinforcement Learning. AGI Research Project,
AGI-2026-001, 2026.

\bibitem{liu2026y2} Z. Liu. Monitor Signal in Cooperative MARL: A
Systematic 4-Pathway Investigation (Lessons-Learned Paper). AGI
Research Project, AGI-2026-001, 2026. \texttt{papers/monitor\_in\_ma\_lessons\_learned.md}

\bibitem{terry2021} J. K. Terry et al. PettingZoo: Gym for
Multi-Agent Reinforcement Learning. NeurIPS 2021.

\bibitem{liu2026hyp} Z. Liu. Y1 9-Hypothesis Framework. AGI Research
Project, AGI-2026-001, 2026. \texttt{papers/y1\_9hypothesis\_framework.md}

\end{thebibliography}

\end{document}
